\documentclass[aspectratio=169,10pt,xcolor={table,dvipsnames}]{beamer}

% UC Berkeley Beamer adaptation copyright 2026 Haichen Wang.
% Based on the Overleaf UC Berkeley Beamer Theme by LianTze Lim and
% Ryan Looney. Licensed under CC BY 4.0. See LICENSE.md.

\usepackage{lmodern}
\usetheme{UCBerkeley}
\usepackage{amsmath,amssymb,booktabs,array,graphicx,tikz}

\title{Scientific presentation title}
\subtitle{Analysis or project subtitle}
\author{Presenter Name}
\institute{Department of Physics, University of California, Berkeley}
\date{Meeting Name\\\today}

\begin{document}

\begin{frame}
\vspace{0.35cm}
{\usebeamerfont{title}\color{Medalist}\inserttitle\par}
\vspace{0.18cm}
\slideopening{This talk presents the scientific question, method, and current result.}
\vspace{0.22cm}
{\large\color{BerkeleyBlue}\insertsubtitle\par}
\vspace{0.30cm}
{\color{BerkeleyBlue}\insertauthor\\\insertinstitute\par}
\vspace{0.25cm}
{\color{BerkeleyBlue}\insertdate\par}
\vspace{0.14cm}
{\scriptsize\color{BerkeleyGrey}UC Berkeley Beamer adaptation by Haichen Wang\par}
\vspace{0.18cm}
\slideclosing{The main result and its current validation status are stated at the start.}
\end{frame}

\begin{frame}{Outline}
\slideopening{The talk is organized around the question, evidence, and conclusion.}
\begin{enumerate}
  \item Scientific question and inputs
  \item Method and supporting evidence
  \item Result, limitations, and next steps
\end{enumerate}
\slideclosing{We begin with the scientific question and the information used to address it.}
\end{frame}

\begin{frame}{Evidence}
\slideopening{The method and the evidence used to evaluate it are summarized here.}
\begin{columns}[T,onlytextwidth]
\column{0.38\textwidth}
\begin{itemize}
  \item State the input or sample.
  \item Define the method or comparison.
  \item Give the validation criterion.
\end{itemize}
\column{0.59\textwidth}
\begin{block}{Primary visual}
\centering
\vspace{1.15cm}
Replace this box with a plot, table, or diagram.
\vspace{1.15cm}
\end{block}
\end{columns}
\slideclosing{State the supported result, limitation, or decision in this sentence.}
\source{Add the public reference or collaboration source needed by the audience.}
\end{frame}

\begin{frame}{Summary and next steps}
\slideopening{The result, its present limitation, and the next action are summarized here.}
\begin{itemize}
  \item Result: state the principal quantitative or technical outcome.
  \item Status: label the result as internal, preliminary, expected, or published.
  \item Next step: name the validation, input, or decision that remains.
\end{itemize}
\begin{alertblock}{Caveat: required check}
State what must be checked and why it affects the interpretation.
\end{alertblock}
\slideclosing{End with the supported takeaway and the next decision or deliverable.}
\end{frame}

\end{document}
